\section{Introduction}\label{intro}

Reading mathematical formulas imposes high cognitive demands: readers must decode symbolic structure, infer relationships between variables, and mentally simulate behavior---imagining how outputs shift as inputs vary or how an algorithm iterates toward a solution~\cite{adams2003reading, Osterholm2006-al}. Interaction can help. Research on learning with multiple representations shows that dynamically linking formulas to visualizations supports comprehension: coordinated representations help learners constrain interpretation and build deeper understanding~\cite{Ainsworth2006-pu}. Direct manipulation can further reduce cognitive effort by closing the gap between a reader's intentions and the system's response, letting users act on mathematical objects and observe outcomes rather than mentally simulating them~\cite{Hutchins1985-wb}.

These findings are reflected in the growing adoption of interactive mathematical content on sites such as Distill.pub and Setosa.io. However, \textit{authoring} these interactions remains prohibitively difficult (Section \ref{FormativeStudy}). Creating a simple interactive formula requires verbose imperative code that entangles computation, state management, interaction logic, and rendering---code that bears no resemblance to the underlying math. To understand this design space, we reviewed 46 interactive formula instances from leading math communication venues including Distill~\cite{distill2017}, Setosa~\cite{setosa2013}, Explorable Explanations~\cite{victor2011explorable}, and others. We found that despite very similar shared types of interactions, nearly all are implemented with custom, case-specific code with no shared abstractions (Section \ref{FormativeStudy}).

Our contribution is \textit{Delta}, a domain-specific language implemented as a React library for authoring interactive mathematical content on the web. Delta is organized around three modes of formula interaction distilled from our analysis of existing interactive explanations: \textit{substitution}, where readers drag symbols directly in a formula and observe how values 
change; \textit{walkthroughs}, where computation \jeff{The introduction of computation here as a thing that happens to formulas is surprising to me. Maybe ``evaluation'' or ``simplification''?} is progressively revealed one step at a time; and \textit{synchronization}, where formula state \jeff{again, ``state'' is not really a thing that formulas have. I think this explanation is mixing ``modes of interaction'' and ``steps to implement interaction''} propagates bidirectionally to linked visualizations. To support these modes, Delta introduces a design that centers on 
five primitives grounded \jeff{$\leftarrow$ What does grounded mean here?}: \textit{formulas} (the \LaTeX{} \jeff{You probably want to say ``notation''?} 
to render), \textit{variables} (which symbols are interactive, their ranges, and display), \textit{semantics} (how values relate computationally), \textit{steps} (how computation is progressively revealed), and \textit{connections} (how formulas synchronize with visualizations). From these primitives, the system derives state management, interactivity, rendering, and synchronized updates. \jeff{I like the punchiness of the preceding two sentences a lot} Beyond the infrastructure, \jeff{own this: ``we also contribute''} an additional contribution is an articulation of what kinds of interactivity should be supported in formulas, distilled from our examples analysis.

We evaluate Delta through two studies. In a controlled comparative study ($n = 19$), participants completed a formula augmentation task faster with Delta than with a baseline React implementation and rated Delta as significantly easier, better suited, and a closer match to how they would describe the task ($p \leq .014$ on all three measures). In an observation study ($n = 10$), participants built interactive formulas from scratch for two mathematical concepts, and we found that Delta's declarative structure is intuitive, scaffolded authoring, and supported pedagogical reasoning. Evaluating the study's data through the framework of the cognitive dimensions of notations, Delta reduced viscosity, diffuseness, and error-proneness \andrew{mismatch between cog dimensions mentioned here and in the results} while exhibiting high closeness of mapping relative to the baseline.

In summary, this paper contributes:
\begin{itemize}
    \item The design and implementation of Delta, a DSL that reifies primitives into a system for authoring where the process of enlivening formulas is made easier and faster.
    \item A system consisting of runtime and React library that supports the authoring of interactive formulas and any linked dynamic representations.
    \item Evidence from usability studies that Delta reduces authoring time and effort compared to a React baseline.
\end{itemize}